% Rapport sur le modèle actuel — projet « Livré ou Remboursé »
% Compilation : xelatex rapport_modele_actuel.tex (deux passes)
\documentclass[11pt,a4paper]{article}

\usepackage{fontspec}
\usepackage{polyglossia}
\setdefaultlanguage{french}
\setmainfont{TeX Gyre Pagella}
\setsansfont{TeX Gyre Heros}
\setmonofont{DejaVu Sans Mono}[Scale=0.85]

\usepackage[margin=2.4cm]{geometry}
\usepackage{xcolor}
\usepackage{booktabs}
\usepackage{tabularx}
\usepackage{array}
\usepackage{siunitx}
\sisetup{output-decimal-marker={,}, group-separator={\,}, group-minimum-digits=4}
\usepackage{pgfplots}
\pgfplotsset{compat=1.17}
\usepackage{tcolorbox}
\usepackage{enumitem}
\setlist{itemsep=2pt, topsep=4pt}
\usepackage{titlesec}
\usepackage{fancyhdr}
\usepackage[hidelinks]{hyperref}
\usepackage{caption}
\captionsetup{font=small, labelfont=bf}

\definecolor{encre}{HTML}{1F3A4D}
\definecolor{sinistre}{HTML}{B5462E}
\definecolor{prudence}{HTML}{C98B1C}
\definecolor{correct}{HTML}{2F6F5E}
\definecolor{gris}{HTML}{6B7680}
\definecolor{fond}{HTML}{EEF2F4}

\titleformat{\section}{\sffamily\Large\bfseries\color{encre}}{\thesection}{0.8em}{}
\titleformat{\subsection}{\sffamily\large\bfseries\color{encre}}{\thesubsection}{0.8em}{}

\pagestyle{fancy}
\fancyhf{}
\fancyhead[L]{\sffamily\small\color{gris}Livré ou Remboursé — modèle actuel}
\fancyhead[R]{\sffamily\small\color{gris}AssurColis}
\fancyfoot[C]{\sffamily\small\thepage}
\renewcommand{\headrulewidth}{0.3pt}

\newtcolorbox{encadre}[1][]{colback=fond, colframe=encre, boxrule=0.5pt, arc=1pt,
  left=8pt, right=8pt, top=6pt, bottom=6pt, fonttitle=\sffamily\bfseries, #1}

\newcommand{\RS}{\,R\$}
\newcommand{\var}[1]{\texttt{#1}}

\begin{document}
\sloppy

% ------------------------------------------------------------------ Titre
\begin{titlepage}
\vspace*{3cm}
{\sffamily\color{gris}\large Master SEP — Data Mining\par}
\vspace{0.6cm}
{\sffamily\bfseries\color{encre}\fontsize{30}{36}\selectfont Livré ou Remboursé\par}
\vspace{0.4cm}
{\raggedright\sffamily\Large Rapport sur le modèle actuel de prédiction des retards de livraison\par}
\vspace{1.2cm}
{\color{encre}\rule{\linewidth}{0.8pt}}
\vspace{0.6cm}

\begin{tabular}{@{}ll}
\textbf{Modèle évalué} & forêt aléatoire \var{foret\_100} (300 arbres) \\
\textbf{Données} & Olist, e-commerce brésilien (Kaggle, CC BY-NC-SA 4.0) \\
\textbf{Entraînement} & commandes de janvier à décembre 2017 (\num{43956}) \\
\textbf{Test} & commandes du 1\textsuperscript{er} janvier au 29 août 2018 (\num{53344}) \\
\textbf{Objet} & état des lieux avant le travail d'amélioration \\
\end{tabular}

\vfill
{\sffamily\color{gris}31 octobre 2026\par}
\end{titlepage}

% ------------------------------------------------------------------ Synthèse
\section*{Synthèse}

\begin{encadre}[title=En une phrase]
Le modèle actuel est \textbf{utile pour l'assureur} (il double le bénéfice par rapport à la
situation sans modèle) mais \textbf{faible statistiquement} (ROC-AUC de 0,641) et
\textbf{sous-estime le risque en 2018}, ce qui laisse une marge d'amélioration nette.
\end{encadre}

\vspace{0.3cm}
\begin{center}
\begin{tabularx}{\linewidth}{@{}lr>{\raggedright\arraybackslash}X@{}}
\toprule
Indicateur (test 2018) & Valeur & Lecture \\
\midrule
Bénéfice de l'assureur & \num{11049}\RS{} & $\times 2{,}1$ par rapport à « proposer à tous » (\num{5191}\RS) \\
Ratio sinistres / primes & 0,68 & 68\RS{} de sinistres pour 100\RS{} de primes \\
Part des commandes à qui on propose & 81,8\,\% & 18,2\,\% de refus \\
ROC-AUC & 0,641 & 0,5 = hasard ; capacité de tri modeste \\
PR-AUC & 0,144 & hasard = 0,087 (taux de sinistre) \\
Score de Brier & 0,0783 & référence sans modèle : 0,0795 \\
Probabilité moyenne prédite & 7,0\,\% & taux réel : 8,7\,\% (sous-estimation) \\
\bottomrule
\end{tabularx}
\end{center}

\paragraph{Ce qui fonctionne.} La démarche est rigoureuse : split temporel, test regardé une
seule fois, choix du modèle sur une métrique métier, contrôle par refus au hasard, contrainte
d'équité respectée dans les cinq régions.

\paragraph{Ce qui limite le modèle.} Une capacité de tri modeste, des probabilités trop basses
dans les tranches risquées, deux variables qui dégradent le classement, et une forte dérive du
taux de retard entre 2017 et 2018.

\paragraph{Premiers essais (section 7).} Retirer les deux variables est neutre et la
calibration dégrade le résultat : le défaut de calibration vient de la dérive temporelle. La
régression logistique trie mieux que la forêt.

\paragraph{Nouveau modèle (section 8).} Une régression logistique avec trois variables
« trajet » (historique des livraisons entre l'État du vendeur et celui du client) porte la
ROC-AUC du test 2018 de 0,641 à \textbf{0,673} et le bénéfice de \num{11049} à
\textbf{12\,238\RS} (+11\,\%). Le plafond de 20\,\% par région est respecté, mais les refus
se concentrent davantage sur certains États (58\,\% à Rio de Janeiro).

\paragraph{Choix d'équité (section 9).} L'équipe a retenu un \textbf{plafond de 20\,\% de refus
par État}. Tous les États sont désormais à 20\,\% de refus au plus ; le bénéfice sur le test
2018 est de \textbf{9\,417\RS} ($\times 1{,}8$ par rapport à « proposer à tous »), soit
\num{2821}\RS{} de moins qu'avec un plafond par région.

\newpage
% ------------------------------------------------------------------ Comparaison
\section*{Modèle initial et modèle final}

\begin{encadre}[title=En une phrase]
Le modèle final \textbf{prédit mieux} et \textbf{traite tous les États de la même façon}, mais
il rapporte \textbf{\num{1632}\RS{} de moins} sur le test 2018 que le modèle initial : ce coût
vient du choix d'équité par État, pas d'une moins bonne qualité de prédiction.
\end{encadre}

\subsection*{Ce qui a changé}

\begin{center}
\begin{tabularx}{\linewidth}{@{}l>{\raggedright\arraybackslash}X>{\raggedright\arraybackslash}X@{}}
\toprule
& Modèle initial & Modèle final \\
\midrule
Algorithme & Forêt aléatoire (300 arbres) & Régression logistique \\
Variables & 11 & 12 : 9 de base + 3 « trajet » \\
Choix du modèle & Bénéfice, 1 fenêtre de validation & ROC-AUC moyenne, 2 fenêtres \\
Règle d'équité & 20\,\% de refus au plus par région & 20\,\% de refus au plus par État \\
Tests automatiques & 18 & 23, dont 4 contre les fuites \\
\bottomrule
\end{tabularx}
\end{center}

\subsection*{Qualité des prédictions}

Ces mesures ne dépendent pas de la règle d'équité.

\begin{center}
\begin{tabular}{@{}lrr@{}}
\toprule
& Initial & Final \\
\midrule
ROC-AUC, validation 2017 (moyenne des 2 fenêtres) & 0,590 & \textbf{0,644} \\
ROC-AUC, test 2018 & 0,641 & \textbf{0,673} \\
PR-AUC, test 2018 (hasard : 0,087) & 0,144 & \textbf{0,150} \\
Score de Brier & 0,0783 & \textbf{0,0776} \\
Probabilité moyenne prédite (réel : 8,7\,\%) & 7,0\,\% & \textbf{9,0\,\%} \\
Taux de retard des 10\,\% les plus risqués & 17,6\,\% & 17,6\,\% \\
\bottomrule
\end{tabular}
\end{center}

\subsection*{Résultats sur le test 2018, chacun avec sa règle}

\begin{center}
\begin{tabular}{@{}lrr@{}}
\toprule
& Initial (par région) & Final (par État) \\
\midrule
Bénéfice & \textbf{11\,049\RS{}} & \num{9417}\RS{} \\
Ratio sinistres / primes & \textbf{0,68} & 0,74 \\
Garantie proposée à & 81,8\,\% & \textbf{83,9\,\%} \\
Refus par État (États d'au moins 200 commandes) & 9 à 39\,\% & \textbf{11 à 20\,\%} \\
Taux de retard des commandes refusées / acceptées & 11,0 / 8,2\,\% & 9,9 / 8,5\,\% \\
\bottomrule
\end{tabular}
\end{center}

\subsection*{À règle égale}

Pour séparer l'effet du modèle de celui de la règle, chaque modèle est évalué avec les deux
règles (bénéfice sur le test 2018) :

\begin{center}
\begin{tabular}{@{}lrrr@{}}
\toprule
& Initial & Final & Écart \\
\midrule
Plafond de 20\,\% par région & \num{11049}\RS{} & \textbf{12\,238\RS{}} & +\num{1189}\RS{} \\
Plafond de 20\,\% par État & \textbf{10\,615\RS{}} & \num{9417}\RS{} & $-$\num{1198}\RS{} \\
\bottomrule
\end{tabular}
\end{center}

\begin{itemize}
  \item \textbf{Le modèle final prédit mieux} sur toutes les mesures, en 2017 comme en 2018, et
        ne sous-estime plus le niveau du risque.
  \item \textbf{Son avantage vient surtout de la géographie.} Avec un plafond par région, il
        gagne \num{1189}\RS{} en refusant davantage à Rio de Janeiro et en Espírito Santo. Le
        plafond par État lui retire cet avantage : à l'intérieur d'un même État, il ne classe pas
        mieux que le modèle initial.
  \item \textbf{Les deux modèles prennent la même décision pour 89\,\% des commandes.}
  \item \textbf{Le modèle final reste celui choisi sur la validation 2017}, où il devance la forêt
        initiale même avec le plafond par État (\num{4188} contre \num{4145}\RS{} de bénéfice
        moyen). Revenir à la forêt parce qu'elle fait mieux en 2018 reviendrait à choisir le modèle
        sur le test.
\end{itemize}

\newpage
% ------------------------------------------------------------------ 1. Problématique
\section{Problématique}

AssurColis, une insurtech fictive, veut proposer aux acheteurs de la plateforme Olist une
\textbf{garantie de livraison} au moment du paiement :
\begin{itemize}
  \item prix unique : \textbf{4\RS} par commande ;
  \item en cas de retard (le \emph{sinistre}), le client est remboursé de ses frais de port et
        reçoit un bon d'achat de \textbf{15\RS}.
\end{itemize}

La question posée au modèle est donc : \emph{au moment du paiement, quelle est la probabilité
que ce colis arrive après la date promise ?} La réponse sert à décider, commande par commande,
si l'on propose la garantie.

\subsection{La règle de décision}

Pour chaque commande, on calcule le coût attendu chargé de la marge de l'assureur :
\[
  \text{coût attendu} = \hat{p}_{\text{retard}} \times (\text{frais de port} + 15) \times 1{,}3
\]
La garantie est proposée si ce coût est inférieur ou égal au prix, 4\RS. Le refus dépend donc
à la fois du risque prédit et du montant de l'indemnité : un colis cher à expédier peut être
refusé même avec un risque moyen.

\subsection{La contrainte d'équité}

L'État du client n'est pas une variable du modèle, mais la distance et le délai promis sont
fortement liés à la géographie. Pour éviter d'exclure de fait les régions éloignées, la règle
impose \textbf{au plus 20\,\% de refus dans chacune des cinq régions du Brésil}. Quand une
région dépasse ce plafond, on re-propose la garantie aux commandes refusées les moins risquées.

\subsection{Hypothèse de souscription}

On suppose que 20\,\% des clients à qui la garantie est proposée l'achètent. Cette hypothèse
n'est pas observable ; elle multiplie tous les bénéfices sans changer le classement des
stratégies.

% ------------------------------------------------------------------ 2. Données
\section{Données}

\subsection{Source et périmètre}

Les données proviennent du jeu public Olist (commandes, clients, vendeurs, produits, paiements,
avis, géolocalisation), chargé dans une base SQLite. Le périmètre retenu va de janvier 2017 à
août 2018, soit \num{97300} commandes livrées dont 7,8\,\% en retard. Avant 2017, les commandes
sont trop rares ; après août 2018, l'issue des commandes n'est pas encore connue.

Le retard a un vrai coût pour la plateforme : \textbf{63\,\% des commandes en retard reçoivent
un avis de 1 ou 2 étoiles}, contre 9\,\% des commandes livrées à temps.

\subsection{Découpage temporel}

\begin{center}
\begin{tabular}{@{}llrr@{}}
\toprule
Rôle & Période & Commandes & Taux de retard \\
\midrule
Choix du modèle : apprentissage & janv.–sept. 2017 & --- & --- \\
Choix du modèle : validation & oct.–déc. 2017 & --- & --- \\
Entraînement final & janv.–déc. 2017 & \num{43956} & 6,8\,\% \\
Test (regardé une seule fois) & janv.–août 2018 & \num{53344} & 8,7\,\% \\
\bottomrule
\end{tabular}
\end{center}

Ce découpage simule une mise en production : le modèle n'a jamais vu de commande postérieure
à celles qu'il doit juger.

\subsection{Variables utilisées}

Seules les informations connues \emph{au moment du paiement} sont utilisées.

\begin{center}
\begin{tabularx}{\linewidth}{@{}l X@{}}
\toprule
Variable & Contenu \\
\midrule
\var{delai\_promis\_jours} & délai annoncé au client entre l'achat et la livraison \\
\var{distance\_km} & distance entre le vendeur et le client \\
\var{meme\_etat} & 1 si le vendeur et le client sont dans le même État \\
\var{frais\_port} & frais de port payés (R\$) \\
\var{prix\_total} & montant des articles (R\$) \\
\var{poids\_total\_g}, \var{volume\_total\_cm3} & poids et volume du colis \\
\var{nb\_articles}, \var{nb\_vendeurs} & complexité de la commande \\
\var{nb\_mensualites} & nombre de paiements échelonnés \\
\var{type\_paiement} & carte, boleto, voucher, débit (catégorielle) \\
\bottomrule
\end{tabularx}
\end{center}

\paragraph{Variables écartées volontairement.} L'État et les coordonnées GPS (équité), le mois
et le jour d'achat (ils apprennent des pics de 2017 qui ne se répètent pas), l'historique de
retard du vendeur et la catégorie de produit (aucun gain en validation).

% ------------------------------------------------------------------ 3. Modèle
\section{Le modèle}

\subsection{Pipeline}

Le modèle est un \var{Pipeline} scikit-learn en deux blocs :
\begin{enumerate}
  \item \textbf{Préparation} (\var{ColumnTransformer}) : pour les dix variables numériques,
        imputation des valeurs manquantes par la médiane puis standardisation ; pour
        \var{type\_paiement}, encodage one-hot. Ces transformations sont apprises sur
        l'entraînement seulement.
  \item \textbf{Estimateur} : \var{RandomForestClassifier} avec 300 arbres,
        \var{min\_samples\_leaf=100}, \var{max\_features="sqrt"}, \var{random\_state=42}.
\end{enumerate}

Le paramètre \var{min\_samples\_leaf=100} impose qu'une feuille regroupe au moins 100 commandes,
ce qui lisse les probabilités et limite le surapprentissage sur une cible rare (7\,\%).

\subsection{Choix du modèle}

Quatre candidats ont été comparés sur la validation d'octobre à décembre 2017, avec comme
critère le \textbf{bénéfice de l'assureur} et non la ROC-AUC.

\begin{center}
\begin{tabular}{@{}lrr@{}}
\toprule
Candidat & ROC-AUC validation & Bénéfice validation \\
\midrule
Régression logistique & \textbf{0,602} & \num{2710}\RS{} \\
Forêt, feuilles $\geq 50$ & 0,569 & \num{2636}\RS{} \\
\textbf{Forêt, feuilles $\geq 100$ (retenue)} & 0,574 & \textbf{2\,742\RS} \\
Forêt, feuilles $\geq 200$ & 0,575 & \num{2608}\RS{} \\
Proposer à tous (sans modèle) & 0,500 & \num{797}\RS{} \\
\bottomrule
\end{tabular}
\end{center}

\begin{encadre}[colframe=prudence, title=Point d'attention]
La forêt retenue devance la régression logistique de seulement \textbf{32\RS} (1,2\,\%) en
bénéfice, alors que la logistique classe mieux (ROC-AUC 0,602 contre 0,574). L'écart est trop
faible pour être significatif sur une validation de trois mois : le choix de la forêt est
fragile.
\end{encadre}

\subsection{Stabilité dans le temps}

Une validation croisée temporelle en quatre plis sur 2017 (on apprend toujours sur le passé et
on juge sur la suite) donne une ROC-AUC de \textbf{0,602 $\pm$ 0,016}
(plis : 0,579 ; 0,611 ; 0,623 ; 0,597). Le modèle est stable, mais à un niveau modeste.

% ------------------------------------------------------------------ 4. Résultats
\section{Résultats sur le test 2018}

\subsection{Comparaison des stratégies}

\begin{center}
\begin{tabular}{@{}lrrr@{}}
\toprule
Stratégie & Proposée à & Bénéfice & S/P \\
\midrule
Proposer à tous (situation actuelle) & 100,0\,\% & \num{5191}\RS{} & 0,88 \\
Règle : refuser au-delà de \num{2000} km & 97,1\,\% & \num{6721}\RS{} & 0,84 \\
Contrôle : refus au hasard (même nombre) & 81,8\,\% & \num{4578}\RS{} & 0,87 \\
Modèle, sans plafond par région & 73,3\,\% & \num{13329}\RS{} & 0,57 \\
\textbf{Modèle, plafond de 20\,\% par région} & \textbf{81,8\,\%} & \textbf{11\,049\RS} & \textbf{0,68} \\
\bottomrule
\end{tabular}
\end{center}

\begin{itemize}
  \item Le modèle bat le contrôle au hasard de \num{6471}\RS{} : son gain vient bien du tri des
        commandes et pas seulement du fait de refuser.
  \item La contrainte d'équité coûte \textbf{2\,279\RS} (\num{13329} contre
        \num{11049}). C'est le prix d'un traitement comparable entre régions.
\end{itemize}

\subsection{Capacité de tri}

\begin{itemize}
  \item \textbf{ROC-AUC 0,641} : en tirant au hasard une commande en retard et une commande à
        l'heure, le modèle donne le risque le plus élevé à la première dans 64\,\% des cas.
  \item \textbf{PR-AUC 0,144}, contre 0,087 pour un tri au hasard : gain réel, mais faible.
  \item Le décile le plus risqué a un taux de retard de 17,6\,\%, cinq fois celui du décile le
        moins risqué (3,5\,\%). Le modèle sépare bien les extrêmes.
\end{itemize}

\subsection{Calibration}

La décision utilise directement la probabilité prédite : si elle est trop basse, le coût
attendu est sous-estimé et l'assureur propose la garantie à des commandes non rentables.

\begin{figure}[h]
\centering
\begin{tikzpicture}
\begin{axis}[
  width=0.82\linewidth, height=6.6cm,
  xlabel={Tranche de risque prédit (décile)}, ylabel={Taux de retard (\%)},
  xmin=0.5, xmax=10.5, ymin=0, ymax=19,
  xtick={1,...,10}, ymajorgrids, grid style={gris!25},
  legend cell align=left, legend style={at={(0.03,0.97)}, anchor=north west, draw=none, fill=white, font=\small},
  tick label style={font=\small}, label style={font=\small},
]
\addplot[ybar, area legend, bar width=12pt, fill=encre!25, draw=encre!60] coordinates {
  (1,3.39) (2,4.38) (3,5.06) (4,5.70) (5,6.30) (6,6.85) (7,7.45) (8,8.18) (9,9.32) (10,13.10)};
\addlegendentry{Probabilité moyenne prédite}
\addplot[mark=*, thick, color=sinistre] coordinates {
  (1,3.51) (2,3.90) (3,5.68) (4,7.16) (5,7.46) (6,8.79) (7,10.10) (8,10.46) (9,12.37) (10,17.64)};
\addlegendentry{Taux de retard observé}
\end{axis}
\end{tikzpicture}
\caption{Calibration sur le test 2018 (\num{53344} commandes, environ \num{5335} par décile).}
\end{figure}

\begin{center}
\begin{tabular}{@{}crrr@{}}
\toprule
Décile & Prédit & Observé & Écart \\
\midrule
1 & 3,4\,\% & 3,5\,\% & +0,1 pt \\
2 & 4,4\,\% & 3,9\,\% & $-$0,5 pt \\
3 & 5,1\,\% & 5,7\,\% & +0,6 pt \\
4 & 5,7\,\% & 7,2\,\% & +1,5 pt \\
5 & 6,3\,\% & 7,5\,\% & +1,2 pt \\
6 & 6,9\,\% & 8,8\,\% & +1,9 pt \\
7 & 7,5\,\% & 10,1\,\% & +2,7 pt \\
8 & 8,2\,\% & 10,5\,\% & +2,3 pt \\
9 & 9,3\,\% & 12,4\,\% & +3,1 pt \\
10 & 13,1\,\% & 17,6\,\% & \textcolor{sinistre}{\textbf{+4,5 pt}} \\
\bottomrule
\end{tabular}
\end{center}

L'ordre des déciles est respecté, mais \textbf{le modèle sous-estime le risque à partir du
quatrième décile}, et d'autant plus que le risque est élevé. En moyenne, il prédit 7,0\,\% de
retards quand 8,7\,\% se produisent. Le score de Brier (0,0783) n'améliore que de 1,5\,\% la
référence sans modèle (0,0795).

\subsection{Importance des variables}

Importance par permutation, mesurée sur le test avec la ROC-AUC (baisse de la ROC-AUC quand on
mélange la variable). Une variable importante est utilisée par le modèle ; ce n'est pas une
cause du retard.

\begin{center}
\begin{tabular}{@{}lrr@{}}
\toprule
Variable & Importance & Écart-type \\
\midrule
\var{distance\_km} & 0,0446 & 0,0016 \\
\var{delai\_promis\_jours} & 0,0440 & 0,0010 \\
\var{meme\_etat} & 0,0340 & 0,0009 \\
\var{frais\_port} & 0,0144 & 0,0009 \\
\var{nb\_articles} & 0,0058 & 0,0004 \\
\var{nb\_vendeurs} & 0,0026 & 0,0005 \\
\var{type\_paiement} & 0,0024 & 0,0005 \\
\var{prix\_total} & 0,0023 & 0,0004 \\
\var{poids\_total\_g} & 0,0005 & 0,0011 \\
\textcolor{sinistre}{\var{nb\_mensualites}} & \textcolor{sinistre}{$-$0,0009} & 0,0002 \\
\textcolor{sinistre}{\var{volume\_total\_cm3}} & \textcolor{sinistre}{$-$0,0014} & 0,0010 \\
\bottomrule
\end{tabular}
\end{center}

Trois variables géographiques (distance, délai promis, même État) portent l'essentiel du
signal. \var{poids\_total\_g} a une importance inférieure à son écart-type, et deux variables
ont une importance \textbf{négative} : les mélanger améliore la ROC-AUC sur 2018. Elles
apportent du bruit appris sur 2017.

\subsection{Équité par région}

\begin{center}
\begin{tabular}{@{}lrrrr@{}}
\toprule
Région & Commandes & Part refusée & Taux de retard & Seuil de coût \\
\midrule
Sudeste & \num{37157} & 17,5\,\% & 7,9\,\% & 4,00\RS{} \\
Sul & \num{7374} & 20,0\,\% & 7,2\,\% & 4,53\RS{} \\
Centro-Oeste & \num{3147} & 20,0\,\% & 8,9\,\% & 5,40\RS{} \\
Norte & 883 & 19,9\,\% & 14,0\,\% & 9,06\RS{} \\
Nordeste & \num{4783} & 20,0\,\% & 16,5\,\% & 10,29\RS{} \\
\bottomrule
\end{tabular}
\end{center}

Le plafond est respecté partout. Il est atteint dans quatre régions sur cinq : seul le Sudeste
reste sous la limite avec la règle de prix normale (seuil de 4\RS). Dans le Nordeste et le
Norte, le seuil est relevé à environ 9 à 10\RS{} : on y propose une garantie à 4\RS{} à des
commandes dont le coût attendu peut atteindre plus du double. C'est là que se concentre le
coût de l'équité.

\subsection{Profils de commandes}

Une segmentation K-means (5 segments, apprise sur 2017) aide à lire les décisions.

\begin{center}
\begin{tabular}{@{}lrrrr@{}}
\toprule
Segment & Part & Retard 2017 & Retard 2018 & Refusées 2018 \\
\midrule
Le voisin (même État) & 30,1\,\% & 4,6\,\% & 5,7\,\% & 2,5\,\% \\
Le petit colis entre États & 36,8\,\% & 7,6\,\% & 11,7\,\% & 10,4\,\% \\
Le grand voyageur (très longue distance) & 14,6\,\% & 6,9\,\% & 7,0\,\% & 11,9\,\% \\
Le colis lourd & 16,3\,\% & 9,2\,\% & 10,8\,\% & 67,7\,\% \\
Le gros panier (plusieurs articles) & 2,3\,\% & 4,7\,\% & 5,8\,\% & 70,9\,\% \\
\bottomrule
\end{tabular}
\end{center}

Les colis lourds et les gros paniers sont refusés à près de 70\,\% alors que leur taux de
retard n'est pas le plus élevé : ils sont refusés surtout parce que leurs frais de port, donc
l'indemnité, sont élevés. À l'inverse, le « petit colis entre États » voit son taux de retard
passer de 7,6\,\% à 11,7\,\% en 2018 mais n'est refusé qu'à 10\,\%.

% ------------------------------------------------------------------ 5. Dérive
\section{Dérive temporelle}

\begin{figure}[h]
\centering
\begin{tikzpicture}
\begin{axis}[
  width=\linewidth, height=6.2cm,
  ybar, bar width=9pt,
  ylabel={Taux de retard (\%)}, ymin=0, ymax=22,
  symbolic x coords={01/17,02/17,03/17,04/17,05/17,06/17,07/17,08/17,09/17,10/17,11/17,12/17,
                     01/18,02/18,03/18,04/18,05/18,06/18,07/18,08/18},
  xtick={01/17,02/17,03/17,04/17,05/17,06/17,07/17,08/17,09/17,10/17,11/17,12/17,01/18,02/18,03/18,04/18,05/18,06/18,07/18,08/18}, x tick label style={rotate=60, anchor=east, font=\scriptsize},
  tick label style={font=\small}, label style={font=\small},
  ymajorgrids, grid style={gris!25},
  enlarge x limits=0.03,
  legend cell align=left, legend style={at={(0.02,0.97)}, anchor=north west, draw=none, font=\small},
]
\addplot[fill=encre!45, draw=none, bar shift=0pt, area legend] coordinates {
  (01/17,5.0) (02/17,4.2) (03/17,6.2) (04/17,8.5) (05/17,4.5) (06/17,4.5) (07/17,4.2)
  (08/17,3.8) (09/17,5.3) (10/17,4.9) (11/17,13.3) (12/17,8.4)};
\addlegendentry{Entraînement (2017)}
\addplot[fill=sinistre!75, draw=none, bar shift=0pt, area legend] coordinates {
  (01/18,6.7) (02/18,14.9) (03/18,20.5) (04/18,5.9) (05/18,7.3) (06/18,1.9) (07/18,4.3) (08/18,6.9)};
\addlegendentry{Test (2018)}
\end{axis}
\end{tikzpicture}
\caption{Taux de retard mensuel, janvier 2017 à août 2018.}
\end{figure}

Le taux de retard n'est pas stable :
\begin{itemize}
  \item en 2017, il oscille entre 4 et 6\,\%, sauf en novembre (13,3\,\%), qui est aussi le mois
        le plus chargé de l'année (\num{7360} commandes, contre environ \num{4200} les mois
        précédents) ;
  \item en 2018, \textbf{février (14,9\,\%) et mars (20,5\,\%)} sont des mois de crise, puis le
        taux redescend jusqu'à 1,9\,\% en juin.
\end{itemize}

Le modèle, appris sur une année plutôt calme, ne peut pas anticiper ces épisodes : les
variables connues au paiement ne disent rien de l'état du réseau de transport à venir. Cette
dérive explique en grande partie la sous-estimation observée en calibration (7,0\,\% prédit
contre 8,7\,\% observé).

% ------------------------------------------------------------------ 6. Bilan
\section{Bilan et pistes d'amélioration}

\subsection{Forces}
\begin{itemize}
  \item Méthode d'évaluation rigoureuse : split temporel, test unique, contrôle au hasard.
  \item Critère de choix aligné sur le métier (bénéfice de l'assureur).
  \item Gain économique réel : \num{11049}\RS, soit $\times 2{,}1$ par rapport à la situation
        actuelle.
  \item Contrainte d'équité respectée et chiffrée (coût de \num{2279}\RS).
  \item Code clair : tous les paramètres sont centralisés dans \var{backend/config.py}.
\end{itemize}

\subsection{Faiblesses}
\begin{itemize}
  \item Capacité de tri modeste (ROC-AUC 0,641).
  \item Probabilités sous-estimées dans les déciles risqués, alors que la décision s'appuie
        directement sur elles.
  \item Deux variables à importance négative, une troisième sans effet mesurable.
  \item Choix forêt contre logistique tranché sur un écart de 32\RS{} seulement.
  \item Sensibilité forte à la dérive temporelle, sans mécanisme de suivi ni de réentraînement.
\end{itemize}

\subsection{Pistes à tester}

Chaque piste sera évaluée de la même façon que le modèle actuel : choix sur la validation 2017,
puis comparaison sur le test 2018 avec le bénéfice de l'assureur comme critère principal.

\begin{center}
\begin{tabularx}{\linewidth}{@{}c X l@{}}
\toprule
Priorité & Piste & Effet attendu \\
\midrule
\textcolor{sinistre}{\textbf{1}} & Retirer \var{volume\_total\_cm3} et \var{nb\_mensualites} & moins de bruit \\
\textcolor{sinistre}{\textbf{1}} & Calibrer les probabilités (\var{CalibratedClassifierCV}, isotonique ou sigmoïde) & coût attendu juste \\
\textcolor{prudence}{\textbf{2}} & Tester \var{HistGradientBoostingClassifier} & meilleur tri \\
\textcolor{prudence}{\textbf{2}} & Ajouter des variables construites (distance / délai promis, frais par kg) & meilleur tri \\
\textcolor{correct}{\textbf{3}} & Revoir le choix forêt / logistique avec plusieurs fenêtres de validation & choix plus robuste \\
\textcolor{correct}{\textbf{3}} & Suivre le taux de retard mensuel et prévoir un réentraînement régulier & résistance à la dérive \\
\bottomrule
\end{tabularx}
\end{center}

Les deux pistes de priorité 1, ainsi que la comparaison forêt / logistique sur plusieurs
fenêtres, ont été testées : la section suivante en rend compte.

% ------------------------------------------------------------------ 7. Premiers essais
\section{Premiers essais d'amélioration}

\subsection{Protocole}

Les essais sont jugés \textbf{uniquement sur 2017} ; le test 2018 n'a pas été relu. Pour ne
pas dépendre d'une seule période, chaque variante est évaluée sur deux fenêtres de
validation, en apprenant toujours sur le passé :

\begin{center}
\begin{tabular}{@{}llrrr@{}}
\toprule
Fenêtre & Apprentissage $\rightarrow$ validation & App. & Val. & Taux de retard \\
\midrule
F1 & janv.–juin 2017 $\rightarrow$ juil.–sept. 2017 & \num{14165} & \num{12350} & 4,4\,\% \\
F2 & janv.–sept. 2017 $\rightarrow$ oct.–déc. 2017 & \num{26515} & \num{17441} & 9,5\,\% \\
\bottomrule
\end{tabular}
\end{center}

La calibration utilise \var{CalibratedClassifierCV} avec des plis temporels
(\var{TimeSeriesSplit}, 3 plis) : chaque calibrateur est appris sur des commandes postérieures
à celles qui ont servi au modèle. Les « 9 variables » sont les 11 variables actuelles sans
\var{volume\_total\_cm3} ni \var{nb\_mensualites}.

\subsection{Résultats}

\begin{center}
\begin{tabular}{@{}lrrrr@{}}
\toprule
& \multicolumn{2}{c}{ROC-AUC} & \multicolumn{2}{c}{Bénéfice} \\
\cmidrule(lr){2-3}\cmidrule(l){4-5}
Variante & F1 & F2 & F1 & F2 \\
\midrule
A. Forêt, 11 variables (actuel) & 0,605 & 0,575 & \num{5711}\RS{} & \num{2750}\RS{} \\
B. Forêt, 9 variables & 0,608 & 0,575 & \num{5710}\RS{} & \num{2729}\RS{} \\
C. Forêt, 9 var. + isotonique & 0,587 & 0,571 & \num{5758}\RS{} & \num{2512}\RS{} \\
D. Forêt, 9 var. + sigmoïde & 0,582 & 0,564 & \num{5779}\RS{} & \num{2448}\RS{} \\
E. Logistique, 11 variables & \textbf{0,630} & \textbf{0,602} & \num{5756}\RS{} & \num{2710}\RS{} \\
F. Logistique, 9 variables & \textbf{0,639} & \textbf{0,601} & \num{5714}\RS{} & \num{2706}\RS{} \\
G. Logistique, 9 var. + isotonique & 0,640 & 0,595 & \num{5763}\RS{} & \num{2529}\RS{} \\
\midrule
Proposer à tous (sans modèle) & 0,500 & 0,500 & \num{5510}\RS{} & 797\RS{} \\
\bottomrule
\end{tabular}
\end{center}

\subsection{Lecture}

\paragraph{Retirer les deux variables ne change rien.} Sur 2017, la ROC-AUC et le bénéfice
sont identiques à quelques unités près. Leur importance négative avait été mesurée sur 2018 ;
s'appuyer sur ce constat pour les retirer reviendrait à régler le modèle sur le test. On peut
les retirer par souci de simplicité, pas pour un gain de performance.

\paragraph{La calibration dégrade le résultat.} Elle fait perdre jusqu'à 300\RS{} sur F2 et
abaisse la ROC-AUC. Le tableau suivant montre pourquoi :

\begin{center}
\begin{tabular}{@{}lrrl@{}}
\toprule
Fenêtre & Proba. moyenne prédite & Taux observé & Lecture \\
\midrule
F1 (juil.–sept.) & 5,6\,\% & 4,4\,\% & le modèle surestime \\
F2 (oct.–déc.) & 5,0\,\% & 9,5\,\% & le modèle sous-estime de moitié \\
\bottomrule
\end{tabular}
\end{center}

\begin{encadre}[colframe=prudence, title=Conclusion principale]
Les probabilités ne sont pas mal formées : c'est le \textbf{niveau général du risque qui change
d'un trimestre à l'autre}, alors que le modèle prédit toujours à peu près le même niveau. Un
calibrateur appris sur le passé ne peut pas deviner le niveau à venir. Le défaut de calibration
observé en 2018 est un symptôme de la dérive temporelle, pas un problème de réglage.
\end{encadre}

\paragraph{La régression logistique trie mieux que la forêt.} Elle gagne 0,025 à 0,03 de
ROC-AUC sur les deux fenêtres, pour un bénéfice équivalent. Le choix initial de la forêt,
tranché sur 32\RS{} d'écart, ne tient pas.

\paragraph{Le bénéfice départage mal les variantes.} Les écarts entre variantes (environ
50\RS{}) relèvent du bruit : avec le plafond de 20\,\% de refus par région, la garantie est
proposée à environ 90\,\% des commandes de validation, ce qui laisse peu de place pour
distinguer les modèles. La ROC-AUC et la PR-AUC sont ici plus informatives.

\subsection{Pistes révisées}

\begin{center}
\begin{tabularx}{\linewidth}{@{}c X l@{}}
\toprule
Priorité & Piste & Statut \\
\midrule
--- & Retirer \var{volume\_total\_cm3} et \var{nb\_mensualites} & neutre, simplification \\
--- & Calibrer les probabilités & écarté \\
\textcolor{sinistre}{\textbf{1}} & Ajouter le taux de retard récent du réseau (30 derniers jours, commandes déjà livrées à la date d'achat) & écarté (section 8) \\
\textcolor{sinistre}{\textbf{1}} & Ajouter la marge du délai promis (délai promis moins délai habituel sur le même trajet, calculé sur le passé) & retenu (section 8) \\
\textcolor{prudence}{\textbf{2}} & Comparer logistique, forêt et \var{HistGradientBoostingClassifier} avec ces variables & logistique retenue \\
\textcolor{correct}{\textbf{3}} & Simuler un réentraînement mensuel sur 2018 & à tester \\
\bottomrule
\end{tabularx}
\end{center}

La variable de taux de retard récent vise directement la cause identifiée : elle donne au
modèle le niveau de risque du moment, qu'il ne voit pas aujourd'hui. Elle exige une précaution
contre la fuite d'information : seules les commandes déjà livrées à la date d'achat peuvent
entrer dans son calcul.

% ------------------------------------------------------------------ 8. Nouveau modèle
\section{Nouveau modèle : les variables « trajet »}

\subsection{Les variables}

Pour chaque commande, on regarde les livraisons passées sur le même trajet (État du vendeur
$\rightarrow$ État du client), telles qu'elles étaient connues \textbf{le 1\textsuperscript{er}
jour du mois d'achat} :
\begin{itemize}
  \item \var{delai\_habituel\_trajet} : délai réel médian entre l'achat et la livraison, sur les
        commandes déjà livrées (médiane globale si le trajet compte moins de 30 livraisons) ;
  \item \var{marge\_delai\_jours} : délai promis au client moins ce délai habituel. Une marge
        faible signifie qu'Olist promet moins que ce que le trajet prend d'habitude ;
  \item \var{taux\_retard\_trajet} : taux de retard passé du trajet, sur les commandes dont la
        date promise est passée, rapproché du taux global pour les trajets peu fréquentés.
\end{itemize}

Une quatrième variable a été testée puis écartée : le \textbf{taux de retard des 30 derniers
jours}. L'issue d'une commande n'est connue qu'à sa date promise, environ trois à quatre
semaines après l'achat : ce signal décrit donc la situation d'il y a un mois. Il indiquait
4,8\,\% en novembre 2017 (13,3\,\% réels) et 20\,\% en avril 2018 (5,9\,\% réels), une fois la
crise passée.

Les deux variables neutres (\var{volume\_total\_cm3}, \var{nb\_mensualites}) ont été retirées :
le modèle compte désormais \textbf{12 variables}. L'absence de fuite est vérifiée par des tests
automatiques (\var{tests/test\_trajets.py}) et par un recalcul à la main sur plusieurs dates.

\subsection{Choix du modèle sur la validation 2017}

Le critère de choix a été fixé avant de regarder le test : la \textbf{ROC-AUC moyenne sur les
deux fenêtres} F1 et F2. Le bénéfice est calculé mais ne départage pas les candidats (section 7).

\begin{center}
\begin{tabular}{@{}lrrrr@{}}
\toprule
& \multicolumn{2}{c}{ROC-AUC} & \multicolumn{2}{c}{Bénéfice} \\
\cmidrule(lr){2-3}\cmidrule(l){4-5}
Candidat (12 variables) & F1 & F2 & F1 & F2 \\
\midrule
\textbf{Régression logistique (retenue)} & \textbf{0,647} & \textbf{0,642} & \num{5819}\RS{} & \num{2658}\RS{} \\
Forêt, feuilles $\geq 100$ & 0,632 & 0,620 & \num{5615}\RS{} & \num{2795}\RS{} \\
Gradient boosting & 0,599 & 0,628 & \num{5614}\RS{} & \num{2959}\RS{} \\
\midrule
\emph{Rappel : forêt actuelle (11 variables)} & 0,605 & 0,575 & \num{5711}\RS{} & \num{2750}\RS{} \\
\bottomrule
\end{tabular}
\end{center}

La validation croisée temporelle sur 2017 donne \textbf{0,645 $\pm$ 0,027}, contre
0,602 $\pm$ 0,016 pour le modèle actuel.

\subsection{Résultats sur le test 2018}

Le test a été regardé une fois, après le choix du modèle.

\begin{center}
\begin{tabular}{@{}lrrr@{}}
\toprule
Indicateur (test 2018) & Modèle actuel & Nouveau modèle & Écart \\
\midrule
ROC-AUC & 0,641 & \textbf{0,673} & +0,032 \\
PR-AUC (hasard : 0,087) & 0,144 & \textbf{0,150} & +0,006 \\
Score de Brier & 0,0783 & \textbf{0,0776} & $-$0,0007 \\
Probabilité moyenne prédite (réel : 8,7\,\%) & 7,0\,\% & 9,0\,\% & \\
\midrule
Bénéfice, plafond de 20\,\% par région & \num{11049}\RS{} & \textbf{12\,238\RS{}} & +\num{1189}\RS{} \\
Ratio sinistres / primes & 0,68 & \textbf{0,64} & $-$0,04 \\
Part des commandes à qui on propose & 81,8\,\% & 80,0\,\% & \\
Bénéfice sans plafond & \num{13329}\RS{} & \num{14428}\RS{} & +\num{1099}\RS{} \\
Coût de l'équité & \num{2279}\RS{} & \num{2189}\RS{} & \\
\bottomrule
\end{tabular}
\end{center}

Pour information, les deux autres candidats ont aussi été évalués sur le test, sans que ce
résultat serve à changer de choix : forêt 0,648 de ROC-AUC et \num{10911}\RS{} ; gradient boosting
0,629 et \num{10576}\RS{}. Le choix fait sur la validation est confirmé.

\paragraph{Calibration.} Le modèle ne sous-estime plus le niveau du risque (9,0\,\% prédit pour
8,7\,\% observé). Il surestime en revanche la tranche la plus risquée (22,0\,\% prédit contre
17,6\,\% observé), ce qui rend l'assureur un peu trop prudent sur ces commandes.

\paragraph{Importance des variables.} Les deux variables trajet les plus utiles arrivent en
tête : \var{marge\_delai\_jours} (0,034) et \var{taux\_retard\_trajet} (0,032), devant le délai
promis (0,023), le fait d'être dans le même État (0,022) et la distance (0,018).

\subsection{Équité : un arbitrage à trancher}

Le plafond est respecté : chaque région a 20\,\% de refus. Il est désormais atteint dans les
cinq régions, Sudeste compris. Sans plafond, le nouveau modèle refuserait 93,5\,\% des commandes
du Nordeste et 87,9\,\% de celles du Norte : il repère encore mieux la géographie que le modèle
actuel.

Surtout, à l'intérieur des régions, les refus se concentrent sur certains États :

\begin{center}
\begin{tabular}{@{}llrrrr@{}}
\toprule
Région & État & Commandes & Retard & Refus actuel & Refus nouveau \\
\midrule
Sudeste & Espírito Santo (ES) & \num{1054} & 14,3\,\% & 30,9\,\% & \textcolor{sinistre}{\textbf{65,6\,\%}} \\
Sudeste & Rio de Janeiro (RJ) & \num{6504} & 15,9\,\% & 25,9\,\% & \textcolor{sinistre}{\textbf{58,5\,\%}} \\
Sudeste & São Paulo (SP) & \num{23485} & 5,6\,\% & 12,5\,\% & 7,6\,\% \\
Nordeste & Alagoas (AL) & 202 & 23,8\,\% & 20,3\,\% & \textcolor{sinistre}{\textbf{52,0\,\%}} \\
Nordeste & Maranhão (MA) & 354 & 21,8\,\% & 34,7\,\% & 48,6\,\% \\
Nordeste & Bahia (BA) & \num{1762} & 15,8\,\% & 9,1\,\% & 9,3\,\% \\
Sul & Santa Catarina (SC) & \num{1892} & 9,2\,\% & 22,3\,\% & 35,8\,\% \\
Sul & Paraná (PR) & \num{2724} & 5,2\,\% & 18,3\,\% & 9,8\,\% \\
\bottomrule
\end{tabular}
\end{center}

Parmi les États d'au moins 200 commandes, l'écart entre le taux de refus le plus haut et le
plus bas passe de 30 à \textbf{58 points}. Les refus suivent mieux le risque réel (Rio de
Janeiro et Espírito Santo ont des taux de retard deux à trois fois supérieurs à São Paulo),
ce qui explique le gain de bénéfice. Mais l'État du client avait été exclu du modèle pour
l'équité, et les variables trajet le réintroduisent indirectement.

\begin{encadre}[colframe=prudence, title=Décision à prendre par l'équipe]
\begin{itemize}[leftmargin=12pt]
  \item \textbf{Garder le nouveau modèle tel quel} et documenter l'écart entre États :
        +\num{1189}\RS{}, plafond par région respecté.
  \item \textbf{Plafonner les refus par État} plutôt que par région : même règle, maille plus
        fine ; le coût est à mesurer sur la validation.
  \item \textbf{Calculer les variables trajet par région} (région du vendeur $\rightarrow$ région
        du client) plutôt que par État : moins précis, moins discriminant entre États.
\end{itemize}
Les deux dernières options doivent être évaluées sur la validation 2017 avant toute nouvelle
lecture du test.
\end{encadre}

% ------------------------------------------------------------------ 9. Choix d'équité
\section{Choix d'équité : un plafond par État}

\subsection{Options évaluées sur la validation 2017}

Deux options ont été comparées au nouveau modèle sur les fenêtres F1 et F2, sans relire le test.
L'écart de refus entre États est mesuré sur les États d'au moins 100 commandes.

\begin{center}
\begin{tabular}{@{}lrrrr@{}}
\toprule
& \multicolumn{2}{c}{ROC-AUC} & \multicolumn{2}{c}{Bénéfice} \\
\cmidrule(lr){2-3}\cmidrule(l){4-5}
Variante & F1 & F2 & F1 & F2 \\
\midrule
Nouveau modèle, plafond par région & 0,647 & 0,642 & \num{5819}\RS{} & \num{2658}\RS{} \\
\textbf{Option 2 : plafond par État} & 0,647 & 0,642 & \num{5743}\RS{} & \num{2633}\RS{} \\
Option 3 : trajets entre régions & 0,641 & \textcolor{sinistre}{0,573} & \num{5732}\RS{} & \textcolor{sinistre}{\num{2315}\RS{}} \\
Options 2 et 3 & 0,641 & 0,573 & \num{5746}\RS{} & \num{2179}\RS{} \\
\midrule
\emph{Forêt actuelle, plafond par État} & 0,607 & 0,574 & \num{5703}\RS{} & \num{2586}\RS{} \\
\bottomrule
\end{tabular}
\end{center}

\begin{itemize}
  \item \textbf{L'option 3 est écartée.} Calculées entre régions, les variables trajet perdent
        leur pouvoir de tri : sur F2, la ROC-AUC tombe à 0,573, sous la logistique sans variable
        trajet (0,601).
  \item \textbf{L'option 2 ne change pas le modèle}, seulement la règle de décision. Sur 2017,
        elle coûte de 25 à 76\RS{} et ramène chaque État à 20\,\% de refus au plus (contre 31
        à 32\,\% au maximum avec le plafond par région). Avec cette règle, le nouveau modèle reste
        devant la forêt actuelle (\num{4188}\RS{} de bénéfice moyen contre \num{4145}\RS{}).
\end{itemize}

\subsection{Résultat sur le test 2018}

\begin{center}
\begin{tabular}{@{}lrrrr@{}}
\toprule
Règle (nouveau modèle) & Proposée à & Bénéfice & S/P & Refus par État \\
\midrule
Sans plafond & 60,8\,\% & \num{14428}\RS{} & 0,44 & --- \\
Plafond de 20\,\% par région & 80,0\,\% & \num{12238}\RS{} & 0,64 & 7,6 à 65,6\,\% \\
\textbf{Plafond de 20\,\% par État (retenu)} & \textbf{83,9\,\%} & \textbf{9\,417\RS{}} & \textbf{0,74} & \textbf{11,2 à 20,0\,\%} \\
\midrule
Proposer à tous & 100,0\,\% & \num{5191}\RS{} & 0,88 & --- \\
\bottomrule
\end{tabular}
\end{center}

\begin{encadre}[colframe=prudence, title=Pourquoi l'écart avec la validation]
Sur 2017, passer de la région à l'État coûtait moins de 80\RS{}. Sur 2018, il coûte
\textbf{2\,821\RS{}}. En 2018, les retards ont fortement augmenté à Rio de Janeiro (15,9\,\%)
et en Espírito Santo (14,3\,\%), contre 5,6\,\% à São Paulo : le gain du nouveau modèle venait
surtout de refus concentrés sur ces États, ce que le plafond par État interdit. Avec la même
règle, l'ancienne forêt aurait obtenu \num{10615}\RS{} sur 2018. Revenir à la forêt pour cette
raison reviendrait à choisir le modèle sur le test : le choix reste celui de la validation.
\end{encadre}

\subsection{Bilan}

\begin{itemize}
  \item La règle est cohérente avec le principe d'équité du projet : l'État du client est exclu
        des variables, et aucun État ne peut être refusé à plus de 20\,\%.
  \item Le coût total de l'équité est de \textbf{5\,011\RS{}} sur 8 mois (\num{14428}
        contre \num{9417}\RS{}). La garantie reste rentable : bénéfice multiplié par 1,8 et
        ratio sinistres / primes de 0,74.
  \item Le seuil de coût est fixé par État : 4,00\RS{} à São Paulo, de 4,7 à 9\RS{} dans le reste
        du Sudeste, le Sul et le Centro-Oeste, de 7,5 à 16\RS{} dans le Norte et de 12,7 à
        28,0\RS{} dans le Nordeste.
  \item À suivre en production : la part de refus par État chaque mois, en particulier dans les
        petits États où elle varie beaucoup.
\end{itemize}

\end{document}
